\documentclass[12pt,usepdftitle=false,aspectratio=1610]{beamer}

\usetheme{Boadilla}
%\usetheme{Madrid}
%\usecolortheme{seahorse}

\setbeamertemplate{navigation symbols}{}

%\usefonttheme[onlymath]{serif}
\usefonttheme[]{serif}
%\useinnertheme{circles}
\setbeamertemplate{caption}[numbered]

\usepackage[T2A]{fontenc}			
\usepackage[utf8]{inputenc}			
\usepackage[english,russian]{babel}
\usepackage{textcomp}
\usepackage{gensymb}
\usepackage{mathtext,cmap,multirow,amsmath,tcolorbox} 				
\usepackage{graphicx,wrapfig,subfig,mathtools,gensymb}
\usepackage[font=small,labelfont=bf]{caption}
\usepackage{minted}

\DeclarePairedDelimiter{\norm}{\lVert}{\rVert} 

\graphicspath{{./figs/09}}

\title[Лекция 9]{Калибровка камеры}

\author{Александр Танченко}
\institute{}
\date{2025}

\begin{document}

%------------------------------------------------------------
\begin{frame}
    \titlepage
\end{frame}

%------------------------------------------------------------
\begin{frame}
    \frametitle{Постановка задачи}
    \begin{figure}
        \centering
        \includegraphics[height=0.3\textheight]{fig10.pdf}
        \includegraphics[height=0.3\textheight]{fig11.pdf}
    \end{figure}
    \begin{itemize}
        \item по соответствиям между 3D точками $\mathbf{X}_i$ и 2D точками $\mathbf{x}_i$ на изображении
        $$
            \mathbf{X}_i 
            \quad\leftrightarrow\quad
            \mathbf{x}_i
        $$
        (точки $\mathbf{X}_i$ не должны лежать в одной плоскости)
        \item построим матрицу камеры $P$
        $$
            \mathbf{x}_i = P\cdot\mathbf{X}_i
        $$
        \item по матрице $P$ найдем внутренние и внешние параметры камеры
        $$
            P\qquad\mapsto\qquad 
            K\,,\,\,
            (R\,,\mathbf{t})
        $$
    \end{itemize}
\end{frame}

%------------------------------------------------------------
\begin{frame}
    \frametitle{Вычисление матрицы камеры P}
    $$
        \mathbf{x}_i = P\cdot\mathbf{X}_i=
        \begin{bmatrix*}[r]
                p_{11} & p_{12} & p_{13} & p_{14}\\
                p_{21} & p_{22} & p_{23} & p_{24}\\
                p_{31} & p_{32} & p_{33} & p_{34}
        \end{bmatrix*}\cdot\mathbf{X}_i
        =
        \begin{bmatrix*}[r]
            \mathbf{P}_1\\ \mathbf{P}_2\\ \mathbf{P}_3 
        \end{bmatrix*}\cdot\mathbf{X}_i
        \,,\qquad
        \mathbf{x}_i=
        \begin{bmatrix*}[r]
            x_i\\ y_i\\ w_i 
        \end{bmatrix*}
    $$
    \medskip

    $$
        \begin{bmatrix*}[r]
            \mathbf{0}^T & -w_i\mathbf{X}_i^T & y_i\mathbf{X}_i^T\\
            w_i\mathbf{X}_i^T & \mathbf{0}^T & -x_i\mathbf{X}_i^T\\
            -y_i\mathbf{X}_i^T & x_i\mathbf{X}_i^T & \mathbf{0}^T
        \end{bmatrix*}
        \begin{bmatrix*}[r]
            \mathbf{P}_1\\ \mathbf{P}_2\\ \mathbf{P}_3 
        \end{bmatrix*}=
        \mathbf{0}
    $$
    \medskip

    $$
        \begin{bmatrix*}[r]
            \mathbf{0}^T & -w_i\mathbf{X}_i^T & y_i\mathbf{X}_i^T\\
            w_i\mathbf{X}_i^T & \mathbf{0}^T & -x_i\mathbf{X}_i^T
        \end{bmatrix*}
        \begin{bmatrix*}[r]
            \mathbf{P}_1\\ \mathbf{P}_2\\ \mathbf{P}_3 
        \end{bmatrix*}=
        \mathbf{0}
    $$
    \medskip

    $$
        A_i\cdot\mathbf{p}=\mathbf{0}
    $$
\end{frame}

%---------------------------------------------------------------------
\begin{frame}
\frametitle{Вычисление матрицы P c помощью DLT алгоритма}
\begin{itemize}
    \item для $n$ соответствий, получим $2n\times 12$ систему однородных линейных уравнений 
    $$
        A\cdot\mathbf{p}=
        \begin{bmatrix}
            A_1\\
            \vdots\\
            A_n
        \end{bmatrix}
        \mathbf{p}=
        \mathbf{0}
    $$
    \item если $n\geqslant 6$, то $\mathrm{rank}(\mathrm{A})=11$
    \item и система будет иметь одно линейно-независимое ненулевое решение $\mathbf{p}$
    \item так как $\mathrm{dof}(P)=11$, то
    \item $\mathbf{p}$ определяется с точностью до множителя
    \item часто нормируют
    $$
        \norm{\mathbf{p}}=1
    $$
\end{itemize}
\end{frame}

%---------------------------------------------------------------------
\begin{frame}
\frametitle{Вычисление матрицы P c помощью DLT алгоритма}
\begin{itemize}
    \item на практике координаты точек определены неточно
    $$
        \mathbf{x}_i=\overline{\mathbf{x}}_i+\varepsilon_i
    $$
    $$
        \mathbf{X}_i=\overline{\mathbf{X}}_i+\varepsilon_i
    $$
    \item ищут приближенное решение системы $A\mathbf{p}=\mathbf{0}$,
    минимизируя алгебраическую ошибку $\norm{A\mathbf{p}}$
    $$
        \mathbf{p}=\arg\min\,
        \norm{A\mathbf{p}}\,,\quad
        \norm{\mathbf{p}}=1
    $$
    \item решением $\mathbf{p}$ будет последний столбец матрицы $V$ 
          в SVD разложении  
          $$
            A=UDV^T
          $$
          (диагональные элементы $D$ упорядочены по убыванию)
\end{itemize}
\end{frame}

%---------------------------------------------------------------------
\begin{frame}
\frametitle{Нормализация данных}
\begin{itemize}
    \item элементы матриц $A_i$ могут принимать как очень маленькие так и огромные значения
    \item перед DLT координаты точек нормализуют c помощью преобразований подобий $S_1$, $S_2$
    $$
        \widetilde{\mathbf{X}}_i=S_1\mathbf{X}_i\,,\qquad
        \widetilde{\mathbf{x}}_i=S_2\mathbf{x}_i
    $$
    так чтобы
    $$
        \mathrm{mean}\{\widetilde{\mathbf{X}}_i\}=0\,,\quad 
        \mathrm{std}\{\widetilde{\mathbf{X}}_i\}=1
    $$
    $$
        \mathrm{mean}\{\widetilde{\mathbf{x}}_i\}=0\,,\quad 
        \mathrm{std}\{\widetilde{\mathbf{x}}_i\}=1
    $$
    \item с помощью DLT вычисляют матрицу камеры $\widetilde{P}$
    \item матрица камеры для исходных координат $\mathbf{x}_i$, $\mathbf{X}_i$ равна 
    $$
        P=S_2^{-1}\widetilde{P}S_1
    $$ 
\end{itemize}
\end{frame}

%---------------------------------------------------------------------
\begin{frame}
\frametitle{Вычисление параметров камеры}
\begin{itemize}
    \item структура матрицы камеры
    $$
        P = K[R\,|\,\mathbf{t}] = [M\,|\,\mathbf{p}_3]\,, \qquad
        M = KR\in\mathbf{R}^{3\times3}\,,\quad
        \mathbf{p}_3=K\mathbf{t}\in\mathbf{R}^3
    $$
    \item $K$ и $R$ можно найти с помощью RQ разложения матрицы $M$
    $$
        M = r\cdot q
    $$
    $r$ -- ортогональная матрица, $q$ -- матрица поворота
    \item зная $K$, можно вычислить $\mathbf{t}$
    $$
        \mathbf{t}=K^{-1}\mathbf{p}_3
    $$
\end{itemize}
\end{frame}

%---------------------------------------------------------------------
\begin{frame}
\frametitle{Калибровка камеры с помощью минимизации геометрической ошибки}
\begin{itemize}
    \item более точную калибровку можно получить, минимизируя геометрическую ошибку
    $$
        K, R, \mathbf{t} = \arg\min\sum_i \norm{\mathbf{x}_i-P\mathbf{X}_i}^2\,,\qquad
        P=K\,[R\,|\,\mathbf{t}]
    $$
    \item эта задача нелинейной оптимизации решается итеративными методами,
    такими, как алгоритм Гаусса-Ньютона, Левенберга-Марквардта, etc.
    \item начальное приближение для $K, R, \mathbf{t}$ можно найти из DLT оценки $P$
\end{itemize}
\end{frame}

%------------------------------------------------------------
\begin{frame}
    \frametitle{Радиальная дисторсия камеры}
    \begin{figure}
        \centering
        \includegraphics[height=0.3\textheight]{fig1.png}
    \end{figure}
    \begin{figure}
        \centering
        \includegraphics[height=0.4\textheight]{fig2.pdf}
    \end{figure}
\end{frame}

%------------------------------------------------------------
\begin{frame}
    \frametitle{Тангенциальная дисторсия камеры}
    \begin{figure}
        \centering
        \includegraphics[height=0.3\textheight]{fig5.png}
    \end{figure}
    \begin{figure}
        \centering
        \includegraphics[height=0.4\textheight]{fig6.png}
    \end{figure}
\end{frame}

%------------------------------------------------------------
\begin{frame}
    \frametitle{Fisheye камера}
    \begin{figure}
        \centering
        \includegraphics[height=0.6\textheight]{fig3.jpg}
        \includegraphics[height=0.6\textheight]{fig4.jpg}
    \end{figure}
\end{frame}

%------------------------------------------------------------
\begin{frame}
    \frametitle{Модель радиальной и тангенциальной дисторсии}
    \begin{figure}
            \centering
            \includegraphics[height=0.25\textheight]{fig7.png}
    \end{figure}
    \begin{itemize}
        \item радиальная дисторсия
        \begin{align*}
            x_d &= x\left(1+k_1r^2+k_2r^4+k_3r^6\right)\,,\qquad r=\sqrt{x^2+y^2}\\
            y_d &= y\left(1+k_1r^2+k_2r^4+k_3r^6\right)
        \end{align*}
        \item тангенциальная дисторсия
        \begin{align*}
            x_d &= x+\left[2p_1xy+p_2\left(r^2+2x^2\right)\right]\\
            y_d &= y+\left[p_1\left(r^2+2y^2\right)+2p_2xy\right]
        \end{align*}
        \item коэффициенты дисторсии
        $$  \left(
            k_1\,,\quad
            k_2\,,\quad
            p_1\,,\quad
            p_2\,,\quad
            k_3
            \right)
        $$
    \end{itemize}
\end{frame}

%------------------------------------------------------------
\begin{frame}
    \frametitle{Коррекция радиальной дисторсии (undistortion)}
    \begin{figure}
        \centering
        \includegraphics[height=0.4\textheight]{fig9.png}
    \end{figure}
    \begin{itemize}
        \item коэффициенты дисторсии $(k_1,k_2,p_1,p_2,k_3)$ находятся из условия:
        проекции прямых в $\mathbf{R}^3$ должны быть прямыми на изображении
        \item зная коэффициенты дисторсии, можно найти undistort rectify map $(m_x, m_y)$
        $$
            x = x_d + m_x\,,\qquad
            y = y_d + m_y
        $$
        
        \item после коррекция дисторсии нужно пересчитать матрицу $K$
    \end{itemize}
\end{frame}

\end{document}